\section{Symmetries in Spin Systems}

A symmetry operation transforms the spatial coordinates, momenta, spins, and
external fields of a system.  It is a symmetry of a Hamiltonian only when the
transformed Hamiltonian describes the same physical system.  For a unitary
operation $G$ this condition is
\begin{equation}
    G\mathcal{H}G^{-1}=\mathcal{H},
    \label{eq:unitary-symmetry-condition}
\end{equation}
whereas time reversal is antiunitary and also complex-conjugates scalar
coefficients.

\paragraph{Source status.}
\textbf{Reference needed.}  The general time-reversal and spatial-inversion
operator rules in this section are included provisionally because no approved
general reference is currently available.  The mirror rules are supported by
the EE237 notes, the isotropic-exchange example by ASTRO, and the magnetic
parity--time example by Zhang \emph{et al.}, as cited locally below.

\subsection{Polar and Axial Vectors}

Position $\mathbf{r}$, momentum $\mathbf{p}$, electric field $\mathbf{E}$, and
charge-current density $\mathbf{J}$ are polar vectors.  Orbital angular momentum
$\mathbf{L}=\mathbf{r}\times\mathbf{p}$, spin $\mathbf{S}$, magnetization
$\mathbf{M}$, and magnetic field $\mathbf{B}$ are axial vectors.  For an
orthogonal spatial transformation represented by $R$, the two vector classes
transform as
\begin{equation}
    \mathbf{v}'=R\mathbf{v},
    \qquad
    \mathbf{a}'=(\det R)R\mathbf{a},
    \label{eq:polar-axial-transformation}
\end{equation}
where $\mathbf{v}$ is polar and $\mathbf{a}$ is axial.  Thus they transform in
the same way under proper rotations ($\det R=+1$) and differently under mirrors
and inversion ($\det R=-1$)~\cite{ee237-week11}.

The most frequently used transformations are summarized below.  The entries
describe the vector attached to the spatial point after that point has been
mapped by the operation.
\begin{center}
\renewcommand{\arraystretch}{1.25}
\small
\begin{tabular}{|p{0.14\linewidth}|p{0.22\linewidth}|p{0.27\linewidth}|p{0.19\linewidth}|}
    \hline
    Operation & $\mathbf{r},\mathbf{p}$ & $\mathbf{S},\mathbf{M}$ &
    $\mathbf{E},\mathbf{B}$ \\
    \hline
    Time reversal $\mathcal{T}$ &
    $\mathbf{r}\to\mathbf{r}$, $\mathbf{p}\to-\mathbf{p}$ &
    $\mathbf{S},\mathbf{M}\to-\mathbf{S},-\mathbf{M}$ &
    $\mathbf{E}\to\mathbf{E}$, $\mathbf{B}\to-\mathbf{B}$ \\
    \hline
    Inversion $\mathcal{P}$ &
    $\mathbf{r},\mathbf{p}\to-\mathbf{r},-\mathbf{p}$ &
    $\mathbf{S},\mathbf{M}\to\mathbf{S},\mathbf{M}$ &
    $\mathbf{E}\to-\mathbf{E}$, $\mathbf{B}\to\mathbf{B}$ \\
    \hline
    Mirror $\mathcal{M}_{\mathbf{n}}$ &
    $\mathbf{v}\to\mathbf{v}-2\mathbf{n}(\mathbf{n}\cdot\mathbf{v})$ &
    $\mathbf{a}\to-\mathbf{a}+2\mathbf{n}(\mathbf{n}\cdot\mathbf{a})$ &
    polar $\mathbf{E}$; axial $\mathbf{B}$ \\
    \hline
\end{tabular}
\end{center}

\subsection{Time-Reversal Symmetry}

Time reversal changes $t\to-t$ and reverses momenta and angular momenta.  Spin
is an intrinsic angular momentum, so
\begin{equation}
    \mathcal{T}\mathbf{S}\mathcal{T}^{-1}=-\mathbf{S},
    \qquad
    \mathcal{T}i\mathcal{T}^{-1}=-i.
    \label{eq:time-reversal-spin}
\end{equation}
For a classical magnetization trajectory, the corresponding transformed
trajectory is
\begin{equation}
    \mathbf{m}_{\mathcal{T}}(t)=-\mathbf{m}(-t).
    \label{eq:time-reversed-magnetization}
\end{equation}
A magnetic state can therefore break time-reversal symmetry even when its
nonmagnetic crystal structure does not.  A fixed applied magnetic field, a
fixed current, damping, or another nonequilibrium drive must not be silently
left unchanged when testing time reversal; these time-odd or irreversible
ingredients generally prevent the driven equation from being
time-reversal-invariant.

\paragraph{Source status.}
\textbf{Reference needed.}  Equations~\eqref{eq:time-reversal-spin} and
\eqref{eq:time-reversed-magnetization}, and the general invariance statement
above, still need an approved reference.

\subsection{Parity or Spatial-Inversion Symmetry}

Here parity means full spatial inversion,
\begin{equation}
    \mathcal{P}:\quad \mathbf{r}\to-\mathbf{r}.
    \label{eq:spatial-inversion}
\end{equation}
Because spin is axial, inversion moves a spin to the inverted site without
reversing its direction.  For a spin texture this gives
\begin{equation}
    \mathbf{S}_{\mathcal{P}}(\mathbf{r})
    =\mathbf{S}(-\mathbf{r}).
    \label{eq:inverted-spin-texture}
\end{equation}
The structure has inversion symmetry only when the inverted sites, chemical
species, interactions, and spin pattern reproduce the original structure.
This is distinct from particle-exchange symmetry.  In non-Hermitian coupled-mode
models the symbol $\mathcal{P}$ may instead denote exchange of two subsystems;
its meaning must therefore be stated explicitly.

\paragraph{Source status.}
\textbf{Reference needed.}  Equations~\eqref{eq:spatial-inversion} and
\eqref{eq:inverted-spin-texture}, including the axial-spin convention under
inversion, still need an approved general reference.

\subsection{Mirror Symmetry}

A mirror with unit normal $\mathbf{n}$ reverses the perpendicular component of
a polar vector.  An axial vector behaves oppositely: its components parallel to
the mirror plane reverse, while its perpendicular component is unchanged
according to Eq.~\eqref{eq:polar-axial-transformation}~\cite{ee237-week11}.
For the $xy$ mirror,
\begin{equation}
    (v_x,v_y,v_z)\to(v_x,v_y,-v_z),
    \qquad
    (S_x,S_y,S_z)\to(-S_x,-S_y,S_z).
    \label{eq:xy-mirror-polar-axial}
\end{equation}
The transformed lattice and every axial spin must both be compared with the
original magnetic structure.  A lattice may therefore possess a mirror plane
that is broken by its magnetic order.

\subsection{Spin-Rotation Symmetry}

Spatial rotations move the lattice, whereas spin rotations act in spin space.
For the isotropic exchange Hamiltonian used in ASTRO,
\begin{equation}
    \mathcal{H}_{\mathrm{ex}}
    =\frac{1}{2}\sum_{i\ne j}J_{ij}\,\mathbf{S}_i\cdot\mathbf{S}_j,
    \label{eq:symmetry-isotropic-exchange}
\end{equation}
a common spin rotation $R_s$ leaves every dot product unchanged:
\begin{equation}
    (R_s\mathbf{S}_i)\cdot(R_s\mathbf{S}_j)
    =\mathbf{S}_i\cdot\mathbf{S}_j.
    \label{eq:global-spin-rotation}
\end{equation}
The exchange term consequently has global spin-rotation symmetry when the
couplings $J_{ij}$ do not depend on spin direction~\cite{astro-code}.  Magnetic
anisotropy, Dzyaloshinskii--Moriya interaction, dipolar coupling, an applied
field, or another spin--orbit-coupled term can reduce this continuous symmetry
to the subset of joint spatial and spin operations that preserves the complete
Hamiltonian.

\subsection{Combined Parity--Time Symmetry in a Two-Sublattice Magnet}

For the two-sublattice order parameters defined in
Eq.~\eqref{eq:a-type-order-parameters}, suppose that inversion exchanges
sublattices $A$ and $B$.  Since inversion does not reverse an axial spin,
whereas time reversal does, their actions are
\begin{equation}
\begin{array}{c|cc}
    & \mathbf{m} & \mathbf{n} \\
    \hline
    \mathcal{P} & \mathbf{m} & -\mathbf{n} \\
    \mathcal{T} & -\mathbf{m} & -\mathbf{n} \\
    \mathcal{P}\mathcal{T} & -\mathbf{m} & \mathbf{n}
\end{array}
    \label{eq:pt-order-parameters}
\end{equation}
Thus a compensated antiferromagnetic pattern can be invariant under the
combined operation even when it is not invariant under $\mathcal{P}$ or
$\mathcal{T}$ separately.  In a different use of parity, Zhang \emph{et al.}
define $\mathcal{P}$ as interchange of two coupled ferrimagnetic layers and
obtain parity--time-symmetric N\'eel-vector dynamics at angular-momentum
compensation~\cite{zhang2024nonhermitian}.  These two meanings should not be
conflated.

\paragraph{Source status.}
\textbf{Reference needed.}  Equation~\eqref{eq:pt-order-parameters} is a direct
derivation from the provisional general transformation rules above.  The cited
paper supports only the separate coupled-layer example.

\subsection{Practical Symmetry Test}

To test a proposed symmetry of a spin model:
\begin{enumerate}
    \item map every position and sublattice label;
    \item transform spin as an axial vector and reverse it additionally under
        time reversal;
    \item transform external fields, currents, and boundary conditions;
    \item rebuild every Hamiltonian or dynamical term with the transformed
        quantities; and
    \item compare the complete transformed model and magnetic pattern with the
        originals.
\end{enumerate}
Checking only the atomic lattice or only the spin arrows is insufficient for a
magnetic symmetry.

\paragraph{Source status.}
\textbf{Reference needed.}  The general five-step test is a provisional
pedagogical synthesis and still needs an approved reference.
